% Results, first part: the protocol paragraph, JetClass-II (the pretraining classes) and JetClass
% (the in-distribution downstream test). Plan: docs/PAPER_OUTLINE.md sections 4.0-4.2; figures F2, F3,
% F4 and table T2 of docs/FIGURE_PLAN.md. Every number is a macro from results_generated.tex
% (experiments/FIGS/make_tables.py); "% CHECK:" marks a claim whose direction must be confirmed against
% the macro values at generation.

\section{Results}\label{sec:results}

Each subsection compares the four levels of the class hierarchy first and the controls second, the
controls against the \Kcoarse-class level whose class count they share. Linear probes on frozen features
carry the comparison on the pretraining classes, with the MLP probe beside every linear result that shows
no difference; fine-tuning carries it on the downstream benchmarks. On JetClass-II we compare label sets
by ratios of $1-\mathrm{AUC}$, and on the downstream tasks by each task's background rejection. The
subsections move outward from the pretraining classes. JetClass-II and JetClass come first, then the two public
benchmarks, top tagging and quark/gluon discrimination; fine-tuning with few training jets, the jet mass
and resonance tagging follow.
% CHECK: drop "and resonance tagging" if section 4.7 and F9 are dropped (samples not staged by 2026-11-12).

\subsection{JetClass-II: the pretraining classes}\label{sec:jetclasstwo}

The first question is which distinctions the frozen features keep on the \Kfull{} classes the networks
were pretrained on, scored by linear probes on the class-balanced samples of section~\ref{sec:evaluation}
(sizes in the caption of figure~\ref{fig:f2} and in table~\ref{tab:probes}). Because this task favors
finer label sets by construction, the claims below rest on comparisons that this bias cannot produce: the
\Kres-class model, whose labels keep every resonant class of the \Kfull-class model, and the controls,
which have as many classes as the \Kcoarse-class level.

A linear probe separates the \Kfull{} classes worse at every coarser level of the hierarchy, and the
largest step is the merge of quark flavor between \Kmid{} and \Kcoarse{} classes
(figure~\ref{fig:f2}). On frozen features of the \Kmid-class model it makes
\MlstCmpJcTwoMidOverFullPr{} than on those of the \Kfull-class model, and on the \Kcoarse-class model
\MlstCmpJcTwoCoarseOverFullPr{}. The step from \Kmid{} to \Kcoarse{} classes costs more than the step
from \Kres{} to \Kmid{} classes, although it merges far fewer classes and removes one kind of
distinction where the other removes three: the number of $b$ or $c$ quarks, gluons against light
quarks, and electrons against muons.
% CHECK: MlstCmpJcTwoCoarseOverMid > MlstCmpJcTwoMidOverRes (today 1.269 against 1.197).
The pretraining loss rewards exactly the distinctions in the label set, so a distinction it stops
rewarding is no longer kept in a form a linear classifier can read, and how much that costs depends on
what the distinction is, not on how many classes the merge removes.

\begin{figure}[t]
\centering
\paperfig{10cm}{F2}
\caption{Linear probes on frozen features against the number of classes in pretraining. The vertical
axis is the ratio of $1-\mathrm{AUC}$ (one minus the area under the receiver operating characteristic
curve, averaged over the \Kfull{} JetClass-II classes, each against the rest) to that of the
\Kfull-class model of the same seed; both axes are logarithmic, lower is better, and the dash-dotted
line marks equality. The dashed line joins the four levels of the class hierarchy; the solid line joins
the same levels after fine-tuning on the largest class-balanced training set. Large markers are means
with bars of one standard deviation over \MlstNRunsWord{} pretraining runs, small points single runs. At
\Kcoarse{} classes, beside the \Kcoarse-class level they are read against, sit the \RandVtwoNPartitions{}
random partitions (\RandVtwoRuns{} run each; the larger marker is their mean and standard deviation) and
the partitions with no, a $b$-quark and a random flavor boundary (\RandFlavRuns{} runs each). The test set holds up to
\JcBalTestPerClass{} jets per class and all jets of the \JcBalNShort{} classes with fewer. This is the
\Kfull-class model's own pretraining task, so finer label sets are favored by construction. For scale,
the area under the curve is \MlstLpJcTwoAucFull{} for the \Kfull-class model and
\MlstLpJcTwoAucInitBn{} for the randomly initialized network with its normalization statistics
recomputed.}
\label{fig:f2}
\end{figure}
% CHECK: pending MlstFtCmpJcTwoCoarseOverFullESix (the solid fine-tuning line); the figure script must
% also draw the three controlled-flavor partitions at 17 (PI decision P11), which the mock-up does not.

The \Kres-class model matches the \Kfull-class model on the resonant classes, under the linear probe
(\MlstCmpJcTwoResResOverFullPr{}) and under the MLP probe (\MlstMlpCmpJcTwoResResOverFullPr{};
table~\ref{tab:mlp}), so the equality is not a limit of what a linear classifier can read. Its cost on
JetClass-II sits entirely in the QCD classes it merges (the bottom block of figure~\ref{fig:f3};
table~\ref{tab:probes}). The merge from \Kfull{} to \Kres{} classes moves no boundary between resonant
classes: it removes only the division of QCD jets by their numbers of $b$, $c$ and $s$ quarks, which
the loss then stops rewarding. Of the discriminants in table~\ref{tab:usecase}, this merge removes only
one, a score against QCD jets split by their heavy-flavor content, as in the QCD subclasses that GN3X
distinguishes~\cite{atlasgnx}. Only an analysis that needs that subdivision pays for the merge.
% CHECK: MlstCmpJcTwoQcdResOverFull > 1 by more than 2 SD (today 1.087 +- 0.006).

At a fixed number of classes, which classes are merged decides the cost. A random partition with
\Kcoarse{} classes makes \MlstCmpJcTwoRandOverCoarsePr{} than the \Kcoarse-class level of the
hierarchy, and \MlstCmpJcTwoRandOverMidPr{} than the \Kmid-class level, which has more than twice as
many classes (figure~\ref{fig:f2}).
% CHECK: MlstCmpJcTwoRandOverMid < 1 by more than 2 SD (today 0.918 +- 0.007); if it straddles 1,
% write "as few errors as the 43-class level".
A random partition keeps the class count and the share of training jets of every \Kcoarse-class group
but not the physics that defines it, and it keeps many of the flavor boundaries the \Kcoarse-class level
removes. The partitions with controlled flavor boundaries, at the same class count, lie between the
\Kcoarse-class level and the random partitions, with the partition that has no flavor boundary nearest
the \Kcoarse-class level.
% CHECK: MlstCmpJcTwoFzeroOverCoarse > MlstCmpJcTwoF{one,oner}OverCoarse > MlstCmpJcTwoRandOverCoarse
% (today 0.928/0.935, 0.794/0.775 and 0.786/0.772, 0.724).
Because the partitions were selected to cut across the \Kcoarse-class groups
(section~\ref{sec:labelsets}), they show that the class count alone does not set the cost on the
pretraining classes, not that any grouping into \Kcoarse{} classes would do as well.

Fine-tuning tests whether the gap survives once every weight adapts to the \Kfull{} classes. After
fine-tuning on the largest class-balanced training set, the \Kcoarse-class model makes
\MlstFtCmpJcTwoCoarseOverFullESixPr{} than the \Kfull-class model, where a linear probe on its frozen
features makes \MlstCmpJcTwoCoarseOverFullPr{} (figure~\ref{fig:f2}, solid line).
% CHECK: pending MlstFtCmpJcTwoCoarseOverFullESix; if it is a null ("equal within"), rewrite the first
% clause as "matches the 188-class model (...)".
Fine-tuning trains on all \Kfull{} classes, so every distinction the coarse labels removed is rewarded
again, and what remains is the cost of starting from features that lack them, at this number of
training jets. The partitions are not fine-tuned on JetClass-II, so this comparison is between levels of
the hierarchy only.

Per class, the largest costs fall on decays that differ from another class only in lepton flavor
(figure~\ref{fig:f3}). From the \Kmid-class level on, a decay with electrons shares its class with the
same decay with muons, and the decays with a $\tau$ lepton that decays to an electron share theirs with
those whose $\tau$ lepton decays to a muon. At that level the class with the largest ratio to the
\Kfull-class model is \MlstPerClassMidTopClass{}, and the classes merged with a lepton-flavor
counterpart are costly as a group (median \MlstPerClassMidLeptonMedianPr{}). They are few, and they stay among the costliest at \Kcoarse{} classes.
% CHECK: MlstPerClassMidTopClass is a lepton-flavor class (X->mumu in the natural-composition probes);
% MlstPerClassMidLeptonN is small; MlstPerClassCoarseTopClass is also a lepton-flavor class.
Such a pair of decays differs in the inputs mainly through the particle-type flags of its leptons, so
once the labels merge the two, nothing in the loss asks the backbone to carry which flag was set.
Because the \Kfull-class model separates these classes from the rest with few errors
(table~\ref{tab:perclass}), a modest absolute loss appears as a large ratio.
% CHECK: the 188-class per-class AUC of the lepton-flavor classes in tab:perclass is high.

Summed over classes, quark flavor carries most of the cost. At \Kcoarse{} classes,
\MlstPerClassCoarseNDoubled{} of the \Kfull{} classes at least double their $1-\mathrm{AUC}$, and the
classes that share their group with a decay differing only in quark flavor account for
\MlstPerClassCoarseQuarkShare{} of the total, taken as the sum over the resonant classes of the
logarithm of each class's ratio. Each of these classes loses less than the lepton-flavor classes, but
together they are most of the resonant classes: most decays in JetClass-II come in variants that differ
only in their quark flavors, and the \Kcoarse-class labels merge every such set of variants.
% CHECK: MlstPerClassCoarseQuarkShare > 50% and > MlstPerClassCoarseLeptonShare; MlstPerClassCoarseQuarkN
% is more than half of the resonant classes; MlstPerClassCoarseQuarkMedian < MlstPerClassCoarseLeptonMedian.

The \Kcoarse-class level alone cannot tell whether a class loses because it is merged with a flavor
counterpart or because of the kind of decay it is, since it merges every such class. The random partitions
separate the two, because each merges some flavor counterparts and keeps others apart. Across the
\RandVtwoNPartitions{} partitions, the classes merged with a counterpart in quark flavor make
\MlstPerClassRandFlavMergedOverApartPr{} at the median than those kept apart from all of theirs; for
lepton flavor the same comparison gives \MlstPerClassRandLeptonMergedOverApartPr{}. The cost follows
the merge rather than the kind of decay: a class loses the separation that a merged label no longer rewards.
% CHECK: MlstPerClassRandFlavMergedOverApart and MlstPerClassRandLeptonMergedOverApart > 1 by more than
% 2 SD; if either is a null, say that the cost follows the kind of decay for that flavor, not the merge.

\begin{figure}[p]
\centering
\paperfig{\linewidth}{F3}
\caption{Where the cost of a coarser label set sits, class by class. Each row is one of the \Kfull{}
JetClass-II classes, grouped by its \Kcoarse-class group (brackets on the left), with the classes of
QCD jets (quarks and gluons that do not come from a resonance decay) in a separate block at the bottom. Each column is a configuration: the
\Kres-, \Kmid- and \Kcoarse-class levels of the hierarchy, the \RandVtwoNPartitions{} random partitions, and the three
partitions with controlled flavor boundaries (none, a $b$-quark boundary, a random flavor boundary). The
color is the class's $1-\mathrm{AUC}$ (one minus the area under the receiver operating characteristic
curve of the class against the rest) under a linear probe on frozen features, divided by the same
quantity for the \Kfull-class model pretrained with that seed and averaged over the runs of the
configuration (\MlstNRunsWord{} per level, \RandFlavRuns{} per partition with a controlled flavor boundary,
\RandVtwoRuns{} per random partition). The scale is logarithmic, light at equality and clipped at a factor of four either way. A dot marks a class that shares its group with a decay differing
from it only in quark or lepton flavor; a dash marks a class with fewer errors than in the \Kfull-class
model; hatching marks the \MlstJcBalNThin{} classes with the fewest test jets. The median over classes
of the standard deviation of the \Kcoarse-class ratio over its runs is
\MlstPerClassCoarseSeedSdMedian{}. Every value is in table~\ref{tab:perclass}.}
\label{fig:f3}
\end{figure}
% CHECK: the dash threshold (0.9 in mockup_figures.figure_f3) and the hatching threshold (1,000 test
% jets) need macros before the caption can state them; MlstJcBalNThin classes are all QCD classes.

Decays that differ only in quark flavor show the cost of the merge most directly
(figure~\ref{fig:f4}). On the pair $X\to b\bar b$ against $X\to c\bar c$, frozen features of the
\Kcoarse-class model give a linear probe \MlstCmpBbCoarseOverFullPr{} of the \Kfull-class model. The
\Kmid-class model, whose labels still keep the two decays apart, makes only
\MlstCmpBbMidOverFullPr{}. At \Kcoarse{} classes the two decays share a class and no class boundary
anywhere separates $b$ from $c$ quarks, so the loss no longer rewards the differences in lifetime and
decay multiplicity between $b$ and $c$ hadrons, which reach the network through the displacement and
particle-type inputs. The \Kmid-class labels separate decays with $b$ quarks, decays with $c$ quarks and
decays with neither within every decay topology, and the pair stays readable.

At a fixed number of classes, a flavor boundary anywhere in the label set keeps quark flavor readable,
even for a pair it does not cut. The three partitions with controlled flavor boundaries have
\Kcoarse{} classes and the same group shares, and \RandFlavRuns{} pretraining runs each, so each comparison below
quotes the value less favorable to a difference. On $X\to YY\to bbqq$ against $X\to YY\to ccqq$, the
partition with a random flavor boundary, which keeps every pair of four-quark decays related by
exchanging $b$ and $c$ quarks in one group, has \MlstCmpBbqqFonerOverFzeroPr{} of the partition with no
flavor boundary. The partition with a $b$-quark boundary, the only one of the three that cuts this pair,
does better still: the partition with a random flavor boundary makes \MlstCmpBbqqFonerOverFonePr{} than
it.

Since all three have the same number of classes, the flavor boundary and not the class count sets
this cost. A boundary between four-quark decays with and without $b$ quarks rewards the network for
telling $b$ from $c$ quarks wherever it draws the line, which likely keeps the inputs that carry that
difference in use for every pair; a boundary on the pair itself adds the separation specific to it.

The output layer can lose a discriminant while the frozen features keep most of what it needs. The
boosted $|V_{cb}|$ measurement uses the scores of $X\to bq$ and $X\to cs$ in the denominator of its
discriminant $D_{bc}$~\cite{vcbboosted}, so these classes must be told apart from $X\to bc$. The
\Kres-class features separate them as well as the \Kfull-class features do, under the linear and the MLP
probe (figure~\ref{fig:f4}c). At \Kmid{} classes, where $X\to bc$ and $X\to bq$ share a class and
$D_{bc}$ can no longer be formed (table~\ref{tab:usecase}), a linear probe makes only
\MlstCmpBcMidOverFullPr{}; at \Kcoarse{} classes it makes \MlstCmpBcCoarseOverFullPr{}. The \Kmid-class
labels still separate decays with a $b$ quark from decays with only a $c$ quark, so the features keep
what separates $X\to bc$ from $X\to cs$ and lose only the number of $b$ quarks that separates it from
$X\to bq$. The \Kcoarse-class labels remove flavor altogether.

\begin{figure}[t]
\centering
\paperfig{\linewidth}{F4}
\caption{Linear probes on pairs of JetClass-II classes against the number of classes in pretraining:
(a) $X\to b\bar b$ against $X\to c\bar c$; (b) $X\to YY\to bbqq$ against $X\to YY\to ccqq$; (c)
$X\to bc$ against $X\to bq$ and $X\to cs$, the classes in the denominator of the discriminant of the
boosted $|V_{cb}|$ measurement. Axes, markers, the hierarchy line and the pairing by seed are as in
figure~\ref{fig:f2}, for a binary linear probe on the pair. At \Kcoarse{} classes sit the
\RandVtwoNPartitions{} random partitions (\RandVtwoRuns{} run each, no mean) and the three partitions with
controlled flavor boundaries (\RandFlavRuns{} runs each). A marker is filled where the label set keeps the pair apart, open where it merges the pair, and
filled on the left only where it merges part of it. In (c), the open square beside the \Kres-class
point is the multilayer perceptron probe, divided by the multilayer perceptron probe of the \Kfull-class
model; the \Kres-class model matches the \Kfull-class model under the linear
(\MlstCmpBcResOverFullPr{}) and the multilayer perceptron probe (\MlstMlpCmpBcResOverFullPr{}). In (b),
the partition with no flavor
boundary makes \MlstCmpBbqqFzeroOverCoarsePr{} than the \Kcoarse-class level, and the random partitions
that merge the pair were all pretrained with one seed and those that keep it apart with the other, so
their difference is confounded with the seed. Every value is in table~\ref{tab:pairs}.}
\label{fig:f4}
\end{figure}

\subsection{JetClass: the in-distribution downstream test}\label{sec:jetclass}

JetClass~\cite{part} holds \MlstFtJcOneNSignal{} classes of signal decays and a background of QCD jets;
its jets carry all \FillNFeatures{} particle inputs the networks were pretrained on and come from a
simulation without pileup. Each
pretrained model is fine-tuned once on a JetClass training subset and compared with the same network
trained from scratch, in the form of ParT's Table~1~\cite{part}: accuracy, the AUC averaged over all
pairs of classes, and the rejection of the background at ParT's working point for each signal class
(table~\ref{tab:jetclass}). The caption of table~\ref{tab:jetclass} gives our training and test sizes,
which differ from those of the published rows.

\providecommand{\TableContinued}[2]{\label{#2}\ifnum\getpagerefnumber{#2}>\getpagerefnumber{#1}\relax{\footnotesize\textit{Table~\ref{#1} (continued)}}\par\smallskip\fi}
\begin{center}
\small
\ifdefined\tablabw\else\newlength{\tablabw}\newlength{\tablabtmp}\newlength{\tablabcol}\newlength{\tablabnat}\newlength{\tablabatom}\fi
\setlength{\tablabw}{\linewidth}\setlength{\tablabnat}{0pt}\setlength{\tablabatom}{0pt}
\setlength{\tablabcol}{1\linewidth}\addtolength{\tablabcol}{-2\tabcolsep}
\setlength{\tablabnat}{0pt}\settowidth{\tablabtmp}{jets}\ifdim\tablabtmp>\tablabnat\setlength{\tablabnat}{\tablabtmp}\fi\settowidth{\tablabtmp}{1 M}\ifdim\tablabtmp>\tablabnat\setlength{\tablabnat}{\tablabtmp}\fi\settowidth{\tablabtmp}{2 M}\ifdim\tablabtmp>\tablabnat\setlength{\tablabnat}{\tablabtmp}\fi\addtolength{\tablabcol}{-\tablabnat}\addtolength{\tablabcol}{-2\tabcolsep}
\setlength{\tablabnat}{0pt}\settowidth{\tablabtmp}{\# params}\ifdim\tablabtmp>\tablabnat\setlength{\tablabnat}{\tablabtmp}\fi\settowidth{\tablabtmp}{2.14\,M}\ifdim\tablabtmp>\tablabnat\setlength{\tablabnat}{\tablabtmp}\fi\settowidth{\tablabtmp}{2.21\,M}\ifdim\tablabtmp>\tablabnat\setlength{\tablabnat}{\tablabtmp}\fi\addtolength{\tablabcol}{-\tablabnat}\addtolength{\tablabcol}{-2\tabcolsep}
\setlength{\tablabnat}{0pt}\settowidth{\tablabtmp}{\textbf{0.8468\,$\pm$\,0.0006}}\ifdim\tablabtmp>\tablabnat\setlength{\tablabnat}{\tablabtmp}\fi\settowidth{\tablabtmp}{\textbf{0.8475\,$\pm$\,0.0009}}\ifdim\tablabtmp>\tablabnat\setlength{\tablabnat}{\tablabtmp}\fi\settowidth{\tablabtmp}{0.8271\,$\pm$\,0.0012}\ifdim\tablabtmp>\tablabnat\setlength{\tablabnat}{\tablabtmp}\fi\addtolength{\tablabcol}{-\tablabnat}\addtolength{\tablabcol}{-2\tabcolsep}
\ifdim\tablabcol<\tablabw\setlength{\tablabw}{\tablabcol}\fi
\setlength{\tablabcol}{1\linewidth}\addtolength{\tablabcol}{-2\tabcolsep}
\setlength{\tablabnat}{0pt}\settowidth{\tablabtmp}{jets}\ifdim\tablabtmp>\tablabnat\setlength{\tablabnat}{\tablabtmp}\fi\settowidth{\tablabtmp}{1 M}\ifdim\tablabtmp>\tablabnat\setlength{\tablabnat}{\tablabtmp}\fi\settowidth{\tablabtmp}{2 M}\ifdim\tablabtmp>\tablabnat\setlength{\tablabnat}{\tablabtmp}\fi\addtolength{\tablabcol}{-\tablabnat}\addtolength{\tablabcol}{-2\tabcolsep}
\setlength{\tablabnat}{0pt}\settowidth{\tablabtmp}{\textbf{0.9856\,$\pm$\,$<$0.0001}}\ifdim\tablabtmp>\tablabnat\setlength{\tablabnat}{\tablabtmp}\fi\settowidth{\tablabtmp}{0.9819\,$\pm$\,0.0002}\ifdim\tablabtmp>\tablabnat\setlength{\tablabnat}{\tablabtmp}\fi\settowidth{\tablabtmp}{0.9847\,$\pm$\,0.0001}\ifdim\tablabtmp>\tablabnat\setlength{\tablabnat}{\tablabtmp}\fi\addtolength{\tablabcol}{-\tablabnat}\addtolength{\tablabcol}{-2\tabcolsep}
\ifdim\tablabcol<\tablabw\setlength{\tablabw}{\tablabcol}\fi
\setlength{\tablabcol}{1\linewidth}\addtolength{\tablabcol}{-2\tabcolsep}
\setlength{\tablabnat}{0pt}\settowidth{\tablabtmp}{jets}\ifdim\tablabtmp>\tablabnat\setlength{\tablabnat}{\tablabtmp}\fi\settowidth{\tablabtmp}{1 M}\ifdim\tablabtmp>\tablabnat\setlength{\tablabnat}{\tablabtmp}\fi\settowidth{\tablabtmp}{2 M}\ifdim\tablabtmp>\tablabnat\setlength{\tablabnat}{\tablabtmp}\fi\addtolength{\tablabcol}{-\tablabnat}\addtolength{\tablabcol}{-2\tabcolsep}
\setlength{\tablabnat}{0pt}\settowidth{\tablabtmp}{\textbf{6{,}480\,$\pm$\,530}}\ifdim\tablabtmp>\tablabnat\setlength{\tablabnat}{\tablabtmp}\fi\settowidth{\tablabtmp}{\textbf{6{,}600\,$\pm$\,330}}\ifdim\tablabtmp>\tablabnat\setlength{\tablabnat}{\tablabtmp}\fi\settowidth{\tablabtmp}{\textbf{6{,}670\,$\pm$\,220}}\ifdim\tablabtmp>\tablabnat\setlength{\tablabnat}{\tablabtmp}\fi\addtolength{\tablabcol}{-\tablabnat}\addtolength{\tablabcol}{-2\tabcolsep}
\setlength{\tablabnat}{0pt}\settowidth{\tablabtmp}{\textbf{2{,}650\,$\pm$\,260}}\ifdim\tablabtmp>\tablabnat\setlength{\tablabnat}{\tablabtmp}\fi\settowidth{\tablabtmp}{\textbf{2{,}820\,$\pm$\,40}}\ifdim\tablabtmp>\tablabnat\setlength{\tablabnat}{\tablabtmp}\fi\settowidth{\tablabtmp}{2{,}610\,$\pm$\,160}\ifdim\tablabtmp>\tablabnat\setlength{\tablabnat}{\tablabtmp}\fi\addtolength{\tablabcol}{-\tablabnat}\addtolength{\tablabcol}{-2\tabcolsep}
\ifdim\tablabcol<\tablabw\setlength{\tablabw}{\tablabcol}\fi
\setlength{\tablabcol}{1\linewidth}\addtolength{\tablabcol}{-2\tabcolsep}
\setlength{\tablabnat}{0pt}\settowidth{\tablabtmp}{jets}\ifdim\tablabtmp>\tablabnat\setlength{\tablabnat}{\tablabtmp}\fi\settowidth{\tablabtmp}{1 M}\ifdim\tablabtmp>\tablabnat\setlength{\tablabnat}{\tablabtmp}\fi\settowidth{\tablabtmp}{2 M}\ifdim\tablabtmp>\tablabnat\setlength{\tablabnat}{\tablabtmp}\fi\addtolength{\tablabcol}{-\tablabnat}\addtolength{\tablabcol}{-2\tabcolsep}
\setlength{\tablabnat}{0pt}\settowidth{\tablabtmp}{\textbf{100.8\,$\pm$\,6.2}}\ifdim\tablabtmp>\tablabnat\setlength{\tablabnat}{\tablabtmp}\fi\settowidth{\tablabtmp}{\textbf{101.8\,$\pm$\,1.5}}\ifdim\tablabtmp>\tablabnat\setlength{\tablabnat}{\tablabtmp}\fi\settowidth{\tablabtmp}{\textbf{103.5\,$\pm$\,1.3}}\ifdim\tablabtmp>\tablabnat\setlength{\tablabnat}{\tablabtmp}\fi\addtolength{\tablabcol}{-\tablabnat}\addtolength{\tablabcol}{-2\tabcolsep}
\setlength{\tablabnat}{0pt}\settowidth{\tablabtmp}{\textbf{1{,}090\,$\pm$\,10}}\ifdim\tablabtmp>\tablabnat\setlength{\tablabnat}{\tablabtmp}\fi\settowidth{\tablabtmp}{\textbf{1{,}150\,$\pm$\,70}}\ifdim\tablabtmp>\tablabnat\setlength{\tablabnat}{\tablabtmp}\fi\settowidth{\tablabtmp}{1{,}030\,$\pm$\,40}\ifdim\tablabtmp>\tablabnat\setlength{\tablabnat}{\tablabtmp}\fi\addtolength{\tablabcol}{-\tablabnat}\addtolength{\tablabcol}{-2\tabcolsep}
\ifdim\tablabcol<\tablabw\setlength{\tablabw}{\tablabcol}\fi
\setlength{\tablabcol}{1\linewidth}\addtolength{\tablabcol}{-2\tabcolsep}
\setlength{\tablabnat}{0pt}\settowidth{\tablabtmp}{jets}\ifdim\tablabtmp>\tablabnat\setlength{\tablabnat}{\tablabtmp}\fi\settowidth{\tablabtmp}{1 M}\ifdim\tablabtmp>\tablabnat\setlength{\tablabnat}{\tablabtmp}\fi\settowidth{\tablabtmp}{2 M}\ifdim\tablabtmp>\tablabnat\setlength{\tablabnat}{\tablabtmp}\fi\addtolength{\tablabcol}{-\tablabnat}\addtolength{\tablabcol}{-2\tabcolsep}
\setlength{\tablabnat}{0pt}\settowidth{\tablabtmp}{\textbf{3{,}060\,$\pm$\,390}}\ifdim\tablabtmp>\tablabnat\setlength{\tablabnat}{\tablabtmp}\fi\settowidth{\tablabtmp}{\textbf{3{,}070\,$\pm$\,330}}\ifdim\tablabtmp>\tablabnat\setlength{\tablabnat}{\tablabtmp}\fi\settowidth{\tablabtmp}{\textbf{3{,}160\,$\pm$\,280}}\ifdim\tablabtmp>\tablabnat\setlength{\tablabnat}{\tablabtmp}\fi\addtolength{\tablabcol}{-\tablabnat}\addtolength{\tablabcol}{-2\tabcolsep}
\setlength{\tablabnat}{0pt}\settowidth{\tablabtmp}{10{,}600\,$\pm$\,2{,}700$^\dagger$}\ifdim\tablabtmp>\tablabnat\setlength{\tablabnat}{\tablabtmp}\fi\settowidth{\tablabtmp}{12{,}900\,$\pm$\,1{,}600$^\dagger$}\ifdim\tablabtmp>\tablabnat\setlength{\tablabnat}{\tablabtmp}\fi\settowidth{\tablabtmp}{14{,}100\,$\pm$\,3{,}500$^\dagger$}\ifdim\tablabtmp>\tablabnat\setlength{\tablabnat}{\tablabtmp}\fi\addtolength{\tablabcol}{-\tablabnat}\addtolength{\tablabcol}{-2\tabcolsep}
\ifdim\tablabcol<\tablabw\setlength{\tablabw}{\tablabcol}\fi
\setlength{\tablabcol}{1\linewidth}\addtolength{\tablabcol}{-2\tabcolsep}
\setlength{\tablabnat}{0pt}\settowidth{\tablabtmp}{jets}\ifdim\tablabtmp>\tablabnat\setlength{\tablabnat}{\tablabtmp}\fi\settowidth{\tablabtmp}{1 M}\ifdim\tablabtmp>\tablabnat\setlength{\tablabnat}{\tablabtmp}\fi\settowidth{\tablabtmp}{2 M}\ifdim\tablabtmp>\tablabnat\setlength{\tablabnat}{\tablabtmp}\fi\addtolength{\tablabcol}{-\tablabnat}\addtolength{\tablabcol}{-2\tabcolsep}
\setlength{\tablabnat}{0pt}\settowidth{\tablabtmp}{\textbf{8{,}590\,$\pm$\,1{,}270}}\ifdim\tablabtmp>\tablabnat\setlength{\tablabnat}{\tablabtmp}\fi\settowidth{\tablabtmp}{\textbf{7{,}800\,$\pm$\,340}}\ifdim\tablabtmp>\tablabnat\setlength{\tablabnat}{\tablabtmp}\fi\settowidth{\tablabtmp}{\textbf{8{,}450\,$\pm$\,210}}\ifdim\tablabtmp>\tablabnat\setlength{\tablabnat}{\tablabtmp}\fi\addtolength{\tablabcol}{-\tablabnat}\addtolength{\tablabcol}{-2\tabcolsep}
\setlength{\tablabnat}{0pt}\settowidth{\tablabtmp}{\textbf{378\,$\pm$\,22}}\ifdim\tablabtmp>\tablabnat\setlength{\tablabnat}{\tablabtmp}\fi\settowidth{\tablabtmp}{\textbf{389\,$\pm$\,9}}\ifdim\tablabtmp>\tablabnat\setlength{\tablabnat}{\tablabtmp}\fi\settowidth{\tablabtmp}{$W\to qq'$}\ifdim\tablabtmp>\tablabnat\setlength{\tablabnat}{\tablabtmp}\fi\addtolength{\tablabcol}{-\tablabnat}\addtolength{\tablabcol}{-2\tabcolsep}
\setlength{\tablabnat}{0pt}\settowidth{\tablabtmp}{\textbf{307\,$\pm$\,7}}\ifdim\tablabtmp>\tablabnat\setlength{\tablabnat}{\tablabtmp}\fi\settowidth{\tablabtmp}{\textbf{308\,$\pm$\,1}}\ifdim\tablabtmp>\tablabnat\setlength{\tablabnat}{\tablabtmp}\fi\settowidth{\tablabtmp}{\textbf{312\,$\pm$\,6}}\ifdim\tablabtmp>\tablabnat\setlength{\tablabnat}{\tablabtmp}\fi\addtolength{\tablabcol}{-\tablabnat}\addtolength{\tablabcol}{-2\tabcolsep}
\ifdim\tablabcol<\tablabw\setlength{\tablabw}{\tablabcol}\fi
\setlength{\tablabnat}{0pt}\settowidth{\tablabtmp}{\mbox{pretrained on \Kfull{} classes,} \mbox{fine-tuned}}\ifdim\tablabtmp>\tablabnat\setlength{\tablabnat}{\tablabtmp}\fi\settowidth{\tablabtmp}{\mbox{pretrained on \Kres{} classes,} \mbox{fine-tuned}}\ifdim\tablabtmp>\tablabnat\setlength{\tablabnat}{\tablabtmp}\fi\settowidth{\tablabtmp}{\mbox{no pretraining,} \mbox{trained from scratch}}\ifdim\tablabtmp>\tablabnat\setlength{\tablabnat}{\tablabtmp}\fi
\settowidth{\tablabtmp}{pretrained on \Kfull{} classes,}\ifdim\tablabtmp>\tablabatom\setlength{\tablabatom}{\tablabtmp}\fi\settowidth{\tablabtmp}{pretrained on \Kres{} classes,}\ifdim\tablabtmp>\tablabatom\setlength{\tablabatom}{\tablabtmp}\fi\settowidth{\tablabtmp}{pretrained on \Kcoarse{} classes,}\ifdim\tablabtmp>\tablabatom\setlength{\tablabatom}{\tablabtmp}\fi
\ifdim\tablabw>7cm\setlength{\tablabw}{7cm}\fi\ifdim\tablabnat<\tablabw\setlength{\tablabw}{\tablabnat}\fi\ifdim\tablabatom>\tablabw\setlength{\tablabw}{\tablabatom}\fi
\begin{minipage}{\linewidth}\centering
\captionof{table}{JetClass: Sophon's network pretrained on JetClass-II with 188, 162, 43 and 17 classes, or not pretrained, fine-tuned on one million JetClass jets, in the form of ParT's Table~1: the number of trainable parameters, accuracy, AUC averaged over all pairs of classes (one against one, ParT's convention; on JetClass's equal class sizes it equals the one-against-rest average, checked), and for each signal class the rejection of QCD jets at 50\% signal efficiency (99\% for $H\to\ell\nu qq'$, 99.5\% for $t\to b\ell\nu$) on $D_S = g_S/(g_S+g_{\mathrm{QCD}})$ of equation~(\ref{eq:discr}). Published rows: trained from scratch on 2 million jets and scored on ParT's \LitPartJcTestMillion-million-jet test set (ParT on 100 million jets reaches an accuracy of 0.861 and an AUC of 0.9877); ours scored on two million (the AUC on 500{,}000 of them); parameters as in table~\ref{tab:benchtop}, ours \MlstNParamsJc{}. Mean $\pm$ standard deviation over the \MlstNRunsWord{} pretraining runs. From scratch: three seeds. Bold: the best value of each column, and every value within one standard deviation of it (the larger of the two values' standard deviations; a configuration with two runs is compared by their mean and standard deviation). $^\dagger$: fewer than $\nmin$ QCD jets pass in at least one run (equation~(\ref{eq:cut})), so that run's rejection has a counting uncertainty above 22\%; these values carry no claim and are never bold; the bold of such a column is chosen among its other values. $\pm\,0$: the three runs pass the same number of QCD jets.}
\label{tab:jetclass}
\par\medskip
\begin{adjustbox}{max width=\linewidth}
\begin{tabular}{>{\raggedright\arraybackslash}p{\tablabw}lcc}
\toprule
 & jets & \# params & accuracy \\
\midrule
ParticleNet~\cite{part} & 2 M & 366\,k & 0.828 \\
ParT~\cite{part} & 2 M & 2.14\,M & 0.836 \\
\midrule
\multicolumn{4}{l}{\textit{this work}} \\
\mbox{no pretraining,} \mbox{trained from scratch} & 1 M & 2.21\,M & 0.8271\,$\pm$\,0.0012 \\
\mbox{pretrained on \Kfull{} classes,} \mbox{fine-tuned} & 1 M & 2.21\,M & \textbf{0.8468\,$\pm$\,0.0006} \\
\mbox{pretrained on \Kres{} classes,} \mbox{fine-tuned} & 1 M & 2.21\,M & \textbf{0.8475\,$\pm$\,0.0009} \\
\mbox{pretrained on \Kmid{} classes,} \mbox{fine-tuned} & 1 M & 2.21\,M & 0.8459\,$\pm$\,0.0002 \\
\mbox{pretrained on \Kcoarse{} classes,} \mbox{fine-tuned} & 1 M & 2.21\,M & 0.8422\,$\pm$\,0.0005 \\
\bottomrule
\end{tabular}
\end{adjustbox}
\end{minipage}
\par\bigskip
\begin{minipage}{\linewidth}\centering
\TableContinued{tab:jetclass}{tab:jetclass:2}
\begin{adjustbox}{max width=\linewidth, max totalheight=0.92\textheight}
\begin{tabular}{>{\raggedright\arraybackslash}p{\tablabw}lc}
\toprule
 & jets & AUC \\
\midrule
ParticleNet~\cite{part} & 2 M & 0.9820 \\
ParT~\cite{part} & 2 M & 0.9834 \\
\midrule
\multicolumn{3}{l}{\textit{this work}} \\
\mbox{no pretraining,} \mbox{trained from scratch} & 1 M & 0.9819\,$\pm$\,0.0002 \\
\mbox{pretrained on \Kfull{} classes,} \mbox{fine-tuned} & 1 M & 0.9854\,$\pm$\,0.0001 \\
\mbox{pretrained on \Kres{} classes,} \mbox{fine-tuned} & 1 M & \textbf{0.9856\,$\pm$\,$<$0.0001} \\
\mbox{pretrained on \Kmid{} classes,} \mbox{fine-tuned} & 1 M & 0.9853\,$\pm$\,0.0001 \\
\mbox{pretrained on \Kcoarse{} classes,} \mbox{fine-tuned} & 1 M & 0.9847\,$\pm$\,0.0001 \\
\bottomrule
\end{tabular}
\end{adjustbox}
\end{minipage}
\par\bigskip
\begin{minipage}{\linewidth}\centering
\TableContinued{tab:jetclass}{tab:jetclass:3}
\begin{adjustbox}{max width=\linewidth, max totalheight=0.92\textheight}
\begin{tabular}{>{\raggedright\arraybackslash}p{\tablabw}lcc}
\toprule
 &  & \multicolumn{2}{c}{rejection of QCD} \\
\cmidrule(lr){3-4}
 & jets & $H\to b\bar b$ & $H\to c\bar c$ \\
 &  & $R_{50}$ & $R_{50}$ \\
\midrule
ParticleNet~\cite{part} & 2 M & 5{,}540 & 1{,}681 \\
ParT~\cite{part} & 2 M & 5{,}587 & 1{,}982 \\
\midrule
\multicolumn{4}{l}{\textit{this work}} \\
\mbox{no pretraining,} \mbox{trained from scratch} & 1 M & 4{,}920\,$\pm$\,190 & 1{,}580\,$\pm$\,60 \\
\mbox{pretrained on \Kfull{} classes,} \mbox{fine-tuned} & 1 M & \textbf{6{,}600\,$\pm$\,330} & \textbf{2{,}650\,$\pm$\,260} \\
\mbox{pretrained on \Kres{} classes,} \mbox{fine-tuned} & 1 M & 6{,}060\,$\pm$\,0 & \textbf{2{,}820\,$\pm$\,40} \\
\mbox{pretrained on \Kmid{} classes,} \mbox{fine-tuned} & 1 M & \textbf{6{,}480\,$\pm$\,530} & 2{,}610\,$\pm$\,160 \\
\mbox{pretrained on \Kcoarse{} classes,} \mbox{fine-tuned} & 1 M & \textbf{6{,}670\,$\pm$\,220} & 2{,}390\,$\pm$\,40 \\
\bottomrule
\end{tabular}
\end{adjustbox}
\end{minipage}
\par\bigskip
\begin{minipage}{\linewidth}\centering
\TableContinued{tab:jetclass}{tab:jetclass:4}
\begin{adjustbox}{max width=\linewidth, max totalheight=0.92\textheight}
\begin{tabular}{>{\raggedright\arraybackslash}p{\tablabw}lcc}
\toprule
 &  & \multicolumn{2}{c}{rejection of QCD} \\
\cmidrule(lr){3-4}
 & jets & $H\to gg$ & $H\to 4q$ \\
 &  & $R_{50}$ & $R_{50}$ \\
\midrule
ParticleNet~\cite{part} & 2 M & 90 & 662 \\
ParT~\cite{part} & 2 M & 93 & 761 \\
\midrule
\multicolumn{4}{l}{\textit{this work}} \\
\mbox{no pretraining,} \mbox{trained from scratch} & 1 M & 88.3\,$\pm$\,3.2 & 670\,$\pm$\,40 \\
\mbox{pretrained on \Kfull{} classes,} \mbox{fine-tuned} & 1 M & \textbf{103.6\,$\pm$\,3.8} & \textbf{1{,}090\,$\pm$\,10} \\
\mbox{pretrained on \Kres{} classes,} \mbox{fine-tuned} & 1 M & \textbf{101.8\,$\pm$\,1.5} & \textbf{1{,}150\,$\pm$\,70} \\
\mbox{pretrained on \Kmid{} classes,} \mbox{fine-tuned} & 1 M & \textbf{103.5\,$\pm$\,1.3} & 1{,}030\,$\pm$\,40 \\
\mbox{pretrained on \Kcoarse{} classes,} \mbox{fine-tuned} & 1 M & \textbf{100.8\,$\pm$\,6.2} & 1{,}050\,$\pm$\,20 \\
\bottomrule
\end{tabular}
\end{adjustbox}
\end{minipage}
\par\bigskip
\begin{minipage}{\linewidth}\centering
\TableContinued{tab:jetclass}{tab:jetclass:5}
\begin{adjustbox}{max width=\linewidth, max totalheight=0.92\textheight}
\begin{tabular}{>{\raggedright\arraybackslash}p{\tablabw}lcc}
\toprule
 &  & \multicolumn{2}{c}{rejection of QCD} \\
\cmidrule(lr){3-4}
 & jets & $H\to \ell\nu qq'$ & $t\to bqq'$ \\
 &  & $R_{99}$ & $R_{50}$ \\
\midrule
ParticleNet~\cite{part} & 2 M & 1{,}654 & 4{,}049 \\
ParT~\cite{part} & 2 M & 1{,}609 & \textbf{6{,}061} \\
\midrule
\multicolumn{4}{l}{\textit{this work}} \\
\mbox{no pretraining,} \mbox{trained from scratch} & 1 M & 1{,}270\,$\pm$\,100 & \textbf{6{,}100\,$\pm$\,600} \\
\mbox{pretrained on \Kfull{} classes,} \mbox{fine-tuned} & 1 M & \textbf{3{,}060\,$\pm$\,390} & 12{,}900\,$\pm$\,1{,}600$^\dagger$ \\
\mbox{pretrained on \Kres{} classes,} \mbox{fine-tuned} & 1 M & \textbf{3{,}070\,$\pm$\,330} & 14{,}100\,$\pm$\,3{,}500$^\dagger$ \\
\mbox{pretrained on \Kmid{} classes,} \mbox{fine-tuned} & 1 M & \textbf{3{,}160\,$\pm$\,280} & 14{,}200\,$\pm$\,2{,}100$^\dagger$ \\
\mbox{pretrained on \Kcoarse{} classes,} \mbox{fine-tuned} & 1 M & 2{,}800\,$\pm$\,200 & 10{,}600\,$\pm$\,2{,}700$^\dagger$ \\
\bottomrule
\end{tabular}
\end{adjustbox}
\end{minipage}
\par\bigskip
\begin{minipage}{\linewidth}\centering
\TableContinued{tab:jetclass}{tab:jetclass:6}
\begin{adjustbox}{max width=\linewidth, max totalheight=0.92\textheight}
\begin{tabular}{>{\raggedright\arraybackslash}p{\tablabw}lccc}
\toprule
 &  & \multicolumn{3}{c}{rejection of QCD} \\
\cmidrule(lr){3-5}
 & jets & $t\to b\ell\nu$ & $W\to qq'$ & $Z\to q\bar q$ \\
 &  & $R_{99.5}$ & $R_{50}$ & $R_{50}$ \\
\midrule
ParticleNet~\cite{part} & 2 M & 4{,}673 & 260 & 215 \\
ParT~\cite{part} & 2 M & 4{,}474 & 307 & 236 \\
\midrule
\multicolumn{5}{l}{\textit{this work}} \\
\mbox{no pretraining,} \mbox{trained from scratch} & 1 M & 3{,}940\,$\pm$\,320 & 272\,$\pm$\,6 & 215\,$\pm$\,8 \\
\mbox{pretrained on \Kfull{} classes,} \mbox{fine-tuned} & 1 M & 9{,}410\,$\pm$\,1{,}510$^\dagger$ & 374\,$\pm$\,8 & \textbf{312\,$\pm$\,6} \\
\mbox{pretrained on \Kres{} classes,} \mbox{fine-tuned} & 1 M & \textbf{7{,}800\,$\pm$\,340} & \textbf{389\,$\pm$\,9} & \textbf{308\,$\pm$\,1} \\
\mbox{pretrained on \Kmid{} classes,} \mbox{fine-tuned} & 1 M & \textbf{8{,}450\,$\pm$\,210} & 376\,$\pm$\,6 & \textbf{307\,$\pm$\,7} \\
\mbox{pretrained on \Kcoarse{} classes,} \mbox{fine-tuned} & 1 M & \textbf{8{,}590\,$\pm$\,1{,}270} & \textbf{378\,$\pm$\,22} & 285\,$\pm$\,11 \\
\bottomrule
\end{tabular}
\end{adjustbox}
\end{minipage}
\end{center}


After fine-tuning, the \Kres- and \Kmid-class models match the \Kfull-class model on JetClass: neither
differs from it in any signal class by the rule of section~\ref{sec:uncertainties}
(table~\ref{tab:jetclass}).
% CHECK: MlstFtJcOneResNRejDiff and MlstFtJcOneMidNRejDiff are zero (none).
The \Kmid-class labels give up the number of $b$ and $c$ quarks, gluons against light quarks and
electrons against muons, and none of these losses shows here. JetClass's signal classes differ from one
another and from the background by their number of prongs and by whether they contain $b$ or $c$
quarks, and the \Kmid-class labels keep both. The largest departure, in the Higgs decay to four
quarks, stays within twice its standard deviation over the runs (table~\ref{tab:jetclassall}).
% CHECK: the largest 43/188 departure is H->4q and does not meet the rule (today 0.885 +- 0.062).

At \Kcoarse{} classes the rejection falls in several signal classes: \MlstFtJcOneCoarseRejDropClasses{}.
The Higgs decay to four quarks loses the most, with \MlstFtJcOneRejHfourqCoarseOverFullPr{}. The loss
is spread over classes rather than concentrated in a class defined by flavor, and $H\to b\bar b$ shows
no drop, but its rejection rests on so few passing background jets that its counting uncertainty is
larger than the drops seen elsewhere (table~\ref{tab:jetclassnbkg}).
% CHECK: MlstFtJcOneCoarseRejDropClasses does not contain H->bb and lists several classes; the H->4q drop
% is the largest of them; the H->bb counting uncertainty (about 31 passing jets) exceeds the H->4q drop.
The drops appear only at the level that removes quark flavor; the \Kmid-class model, whose labels still
separate decays with $b$ quarks, with $c$ quarks and with neither, shows none of them. The fine-tuned
\Kcoarse-class network therefore likely starts from features that carry less of the charm and bottom
content that $W$, $Z$ and Higgs decays to quarks often have and that QCD jets rarely have.

On JetClass the random partitions recover part of the gap between the \Kcoarse- and \Kfull-class
models, and how much depends on the number of prongs. With \RandVtwoRuns{} run per partition, we read
these values descriptively. The partitions' fine-tuned accuracy exceeds that of the \Kcoarse-class model by
\MlstFtJcOneAccRandMinusCoarse{} (absolute), which places them between the \Kcoarse- and \Kmid-class
models, whereas on the pretraining classes they did better than the \Kmid-class level. On the two-prong
decays $H\to c\bar c$, $W\to qq'$ and $Z\to q\bar q$ they come close to the \Kfull-class model;
in the Higgs decay to four quarks and the two top-quark classes they stay near the \Kcoarse-class model
(table~\ref{tab:jetclassall}).
% CHECK: today accuracy 188 0.8491, partitions 0.8461, 43 0.8480, 17 0.8434 (curves/full_best70_leg2.json);
% the partitions recover most of the 17-to-188 rejection gap on H->cc, W, Z and little on H->4q and top.
The partitions were built to cut across the \Kcoarse-class groups, and in doing so they put decays with
different numbers of prongs into one group (table~\ref{tab:randprong}). They likely keep the flavor
boundaries that help separate two-prong decays from the background and lose the prong-count boundaries
that the three- and four-prong classes rely on.
